\usepackage[a4paper,
            left=3.57cm,right=2.77cm,
            top=2.75cm,bottom=2.54cm,
            footskip=7.9mm,
            headheight=5mm,headsep=7.5mm]{geometry}
% \usepackage{showframe}
\usepackage{fancyhdr}
\usepackage{titletoc} % 用于设置目录样式的包
\usepackage{graphicx} % 支持插图处理
\usepackage{pdfpages}
\usepackage{pifont}
\usepackage{booktabs} % 用于绘制三线表的横线
\usepackage{amsmath} % 基础数学包
\usepackage{amsfonts} % 支持公式中使用特殊字体
% \usepackage{mathtools} % 基础数学包基础上增加额外命令
% \usepackage{tabularx} % 用于绘制表格
% \usepackage{multirow} % 支持在表格中横跨多列
% \usepackage[hidelinks]{hyperref} % 将交叉引用设置为超链接
\usepackage{newtxmath}
\usepackage{caption}
\usepackage[style=gb7714-2015]{biblatex}

\setmainfont{Times New Roman} 
\setsansfont{SimSun}[AutoFakeBold={3.17}]

\addbibresource{reference.bib}

\DeclareFieldFormat{labelnumberwidth}{
  {\sffamily\mkbibbrackets{#1}}%  使用自定义的无衬线字体命令
}

% 定义一个新的分隔符：使用 \quad（1em 空格）
\DeclareCaptionLabelSeparator{quadspace}{\quad}

% 设置图注：字体为 small，分隔符用刚刚定义的 quadspace
\captionsetup[figure]{
  font=small,
  labelsep=quadspace
}

% 同样设置表格注
\captionsetup[table]{
  font=small,
  labelsep=quadspace
}